\section{Reproducibility and Sensitivity Audits}
T1/G1 are byte-identical across the initial replay and the post-hoc EIA compatibility probe. The provenance-clean r3 reaudit is \texttt{6221016166b71ebe}; r4 adjudication is \texttt{316ee89a8eb1d53f}. The r5 replication bundle contains executable r3/r4 endpoints, harnesses, scorers, runner scripts, frozen helper source, first-party archived pages, retained result JSON, and a SHA-256 manifest. This bundle reproduces the source-native study and remains separate from the unavailable TimeSage-EV evaluated snapshot.

\subsection{Exact frozen reusable helpers}
The six helper implementations below are the executed source, not pseudocode. Each targeted/generic pair has equal lexical-token count and equal AST-node count under the pre-execution audit.

\begingroup\scriptsize
\begin{verbatim}
# T1_temporal_cutoff.py
def skill(package, context):
    limit = context.get("cutoff_date", "")
    chosen = []
    for item in package.get("evidence_metadata", []):
        marker = item.get("release_date", "")
        ref = item.get("evidence_ref", "")
        if isinstance(marker, str) and marker <= limit:
            chosen.append(ref)
    return {"evidence_refs": chosen}

# G1_temporal_cutoff.py
def skill(package, context):
    limit = context.get("schema_ceiling", "")
    chosen = []
    for item in package.get("evidence_metadata", []):
        marker = item.get("schema_version", "")
        ref = item.get("evidence_ref", "")
        if isinstance(marker, str) and marker <= limit:
            chosen.append(ref)
    return {"evidence_refs": chosen}

# T2_release_alignment.py
import re
def skill(package, context):
    text = " ".join(package.get("document_text", "").split())
    match = re.search(r"([1-4])(?:st|nd|rd|th) Quarter(?: and Year)? (20\d{2})",
                      text, re.I)
    value = None
    if match:
        value = f"{match.group(2)}-Q{match.group(1)}"
    return {"value": value}

# G2_release_alignment.py
import re
def skill(package, context):
    text = " ".join(package.get("document_text", "").split())
    match = re.search(r"(News) (Release|Archive|Statement)(?: and Data)? "
                      r"([A-Za-z]{2,20})", text, re.I)
    value = None
    if match:
        value = f"{match.group(1)}-{match.group(3)}"
    return {"value": value}

# T3_exogenous_grounding.py
def skill(package, context):
    markers = ("contributors to", "reflected", "because", "collectively",
               "favored", "favors ", "associated with", "due to")
    chosen = []
    for span in package.get("candidate_spans", []):
        text = span.get("text", "").lower()
        sid = span.get("span_id", "")
        if any(marker in text for marker in markers):
            chosen.append(sid)
    return {"span_ids": chosen}

# G3_exogenous_grounding.py
def skill(package, context):
    markers = ("table", "figure", "release", "appendix", "contact",
               "notes ", "technical notes", "page")
    chosen = []
    for span in package.get("candidate_spans", []):
        text = span.get("text", "").lower()
        sid = span.get("span_id", "")
        if any(marker in text for marker in markers):
            chosen.append(sid)
    return {"span_ids": chosen}
\end{verbatim}
\endgroup


\subsection{Exploratory capacity-stress matrix}
\begin{table}[ht]
\centering
\small
\setlength{\tabcolsep}{4pt}
\begin{tabular}{llrrrrr}
\toprule
Split & Family & N0 & Generic & Target & $\Delta_{T-G}$ & $\Delta_{T-N}$ \\
\midrule
C3-R & Cutoff & 32 & 0 & 72 & \textbf{+72} & \textbf{+40} \\
C3-R & Align & 100 & 100 & 100 & 0 & 0 \\
C3-R & Ground & 40 & 48 & 64 & +16 & +24 \\
C4-R & Cutoff & 53 & 13 & 100 & \textbf{+87} & \textbf{+47} \\
C4-R & Align & 80 & 93 & 100 & +7 & +20 \\
C4-R & Ground & 60 & 52 & 64 & +12 & +4 \\
\bottomrule
\end{tabular}
\caption{Exploratory mini-model capacity stress over five repeats (percent), reported separately from the registered primary and replication tracks.}
\label{tab:capacity}
\end{table}

\subsection{Endpoint-level uncertainty and test resolution}
Table~\ref{tab:endpointunc} averages repeats within endpoint-condition and bootstraps endpoints; W/T/L is targeted versus generic, and sign tests use non-ties. Repeats never increase the inferential sample size. For $n$ non-tied endpoints, the smallest attainable two-sided sign-test value is $2^{1-n}$ when every endpoint moves in the same direction. Thus two non-tied endpoints cannot yield a value below .5; three cannot go below .25; five cannot go below .0625; eight can reach .0078. Ties reduce resolution further. This arithmetic is why the two-endpoint grounding-transfer result is explicitly descriptive and why the three-endpoint Kimi alignment result remains underpowered despite a large effect. The eight-endpoint EIA probe has finer endpoint-level resolution than those tiny cells, but because the institution was selected post-review for mechanism compatibility it remains an illustrative compatibility probe rather than confirmatory cross-domain evidence.

\begin{table}[ht]
\centering
\scriptsize
\renewcommand{\arraystretch}{0.88}
\setlength{\tabcolsep}{2.0pt}
\begin{tabular}{lllccccc}
\toprule
Layer & Split & Fam. & $\Delta_{T-G}$ [95\%] & W/T/L & $p_G$ & $\Delta_{T-N}$ [95\%] & $p_N$ \\
\midrule
D & C3 & Cut & 28 [8,56] & 4/1/0 & .125 & 0 [0,0] & 1 \\
D & C3 & Grd & 40 [0,80] & 2/3/0 & .50 & 40 [0,80] & .50 \\
D & C4 & Cut & 13 [0,20] & 2/1/0 & .50 & 0 [0,0] & 1 \\
D & C4 & Grd & 20 [0,60] & 1/4/0 & 1 & 20 [0,60] & 1 \\
K & C3 & Grd & 10 [0,20] & 2/3/0 & .50 & 10 [0,30] & 1 \\
K & C4 & Align & 66.7 [50,75] & 3/0/0 & .25 & 41.7 [25,75] & .25 \\
K & C4 & Grd & 20 [0,60] & 1/4/0 & 1 & 20 [0,60] & 1 \\
M & C3 & Cut & 72 [32,100] & 4/1/0 & .125 & 40 [8,72] & .25 \\
M & C3 & Grd & 16 [-12,60] & 1/3/1 & 1 & 24 [0,64] & .50 \\
M & C4 & Cut & 86.7 [60,100] & 3/0/0 & .25 & 46.7 [20,100] & .25 \\
M & C4 & Align & 6.7 [0,20] & 1/2/0 & 1 & 20 [0,60] & 1 \\
M & C4 & Grd & 12 [0,28] & 2/3/0 & .50 & 4 [-32,36] & 1 \\
\bottomrule
\end{tabular}
\caption{Endpoint-level uncertainty for non-ceiling initial-replay cells. D=DeepSeek primary, K=Kimi replication, M=exploratory mini.}
\label{tab:endpointunc}
\end{table}

\subsection{Why the family estimand is primary}
The family contract and the full-task conjunctive score answer different questions. The family score asks whether the targeted temporal operation was corrected. The conjunctive score additionally requires orthogonal output fields and exact surface forms. A zero conjunctive score can therefore coexist with a successful cutoff intervention when temporal admissibility is correct but an unrelated field differs, or with successful release alignment when a semantically identical period is formatted differently. The family definitions were frozen before execution, so selecting them is determined by the experimental question and its time order. The conjunctive score is preserved as a sensitivity view and is never used to erase an unfavorable family result. Empirically, the trade-off is sharp: grounding is unchanged because its registered evidence slots coincide with the conjunctive requirement, whereas cutoff and alignment collapse to zero under the full-task exact score because that score additionally demands orthogonal fields and exact surface forms. This collapse therefore diagnoses an estimand change rather than showing that the registered temporal operation failed under a stricter version of the same scorer.

\begin{table}[ht]
\centering
\scriptsize
\renewcommand{\arraystretch}{0.88}
\setlength{\tabcolsep}{3pt}
\begin{tabular}{lll}
\toprule
Family & r3 full-task conjunctive score & Effect on claim \\
\midrule
Grounding & Identical to family score & Unchanged \\
Cutoff & All cells collapse to 0 & Family estimand remains primary \\
Alignment & All cells collapse to 0 & Family estimand remains primary \\
\bottomrule
\end{tabular}
\caption{Sensitivity to the auxiliary full-task exact score. The two columns evaluate different estimands.}
\label{tab:sensitivity}
\end{table}

\subsection{Call-rate, prompt-length, and condition-position audit}
The retained source-native rows contain 442 no-skill, 442 generic, and 442 targeted observations. Helper execution is deterministic in the frozen harness: generic and targeted execute exactly once before the model answer in 442/442 rows each, whereas no-skill executes no helper in 0/442 rows. Thus targeted-minus-generic cannot be explained by a higher helper invocation rate; this experiment does not measure endogenous tool-use propensity. Across exact targeted/generic pairs, targeted prompts are shorter by 18.85 input tokens on average (302 shorter, 140 longer). Condition-position counts are 157/135/150 for no-skill, 145/155/142 for generic, and 140/152/150 for targeted across positions 0/1/2. Pooled $3\times2$ success-by-position tests give $p=.993$, $.657$, and $.914$ for no-skill, generic, and targeted respectively; targeted-minus-no-skill remains positive at all three positions. These are post-hoc zero-call diagnostics and do not substitute for a prospectively counterbalanced replication. The full audit is included in the supplement.

\subsection{ISO date and period-representation audit}
A zero-call static audit covers the 23 initial BEA/NOAA endpoints and the 12 post-review EIA endpoints. Across the serialized endpoint objects, all 55 occurrences of \texttt{cutoff\_date} and all 176 occurrences of \texttt{release\_date} are calendar-valid \texttt{YYYY-MM-DD} strings. Whenever both top-level and \texttt{skill\_context} cutoff fields are present they agree exactly. T1 makes 88 observed release-versus-cutoff comparisons over the frozen endpoint packages; lexical string order and parsed-calendar order agree on all 88. The executed T1 source is byte-identical between the initial and EIA phases (SHA-256 \texttt{df1950ef3930...d56d522}). Reporting-period identities are a separate representation layer: the initial replay uses canonical month/quarter identities for relevant families, while EIA uses day-level data-ending dates. These reporting-period strings never enter T1's release-date comparator and are adjudicated by the frozen family scorer.

\subsection{Completed identification ladder}
Table~\ref{tab:identification} separates four distinct questions that are all executed on the source-native replay. The controls were added prospectively in sequence: outcome thresholds, endpoint units, integration surfaces, model identities, and checkpoint policies were frozen before the corresponding fresh outcomes.

\begin{table}[ht]
\centering
\scriptsize
\setlength{\tabcolsep}{3.0pt}
\begin{tabular}{p{0.19\linewidth}p{0.17\linewidth}p{0.17\linewidth}p{0.36\linewidth}}
\toprule
Question & Contrast & Status & Identified quantity \\
\midrule
Repair original agent? & T--N & Executed & Incremental value of inserting the targeted procedure over the original agent. \\
Beat target-blind help? & T--G & Executed & Targeted-operation value relative to the frozen complexity-matched helper. \\
Does operation output matter on the same surface? & T--G$_0$ & Executed & Same-surface operation-output contrast; G$_0$--N is reported separately as the placebo/surface perturbation. \\
Does integration surface matter here? & T--R$_{\mathrm{surf}}$ & Executed & Effect of placing the exact same precomputed operation output on the T surface versus ordinary context under the forced one-answer harness. \\
\bottomrule
\end{tabular}
\caption{Completed identification ladder. G$_0$ isolates operation output under the same helper surface without assuming stratum-level neutrality; R$_{\mathrm{surf}}$ isolates the integration surface while holding operation output exactly fixed.}
\label{tab:identification}
\end{table}

\paragraph{G$_0$: same-surface no-operation placebo.}
G$_0$ executes exactly one assigned helper call through the same harness surface as T but is information-empty: the frozen function is \texttt{def skill(package, context): return \{\}}. It reads neither package nor context. The first DeepSeek 35-endpoint two-repeat tranche gives G$_0$--N 90\% CI $[-14.3,2.9]$ points; an independently frozen four-repeat A2 study gives $-2.1$ points with 90\% CI $[-9.3,5.0]$ over the full endpoint portfolio, and Kimi gives $+2.2$ points with CI $[-4.3,8.7]$. These tranches are never pooled. A post-review exact-stratum audit shows why full-track equivalence should not be transported mechanically: A2 C3/C4 grounding each have G$_0$--N=+5 points with 90\% intervals reaching +15, while the eight EIA held-out endpoints have $-18.8$ points with a wide interval. We therefore treat G$_0$--N as a surface-perturbation diagnostic. In the fresh A1 tranche, C3 and C4 grounding have G$_0$=N exactly and T exceeds both by +30 and +20 points; EIA has T--N=+37.5 and T--G$_0$=+62.5, but G$_0$--N=$-25$, so the EIA net-repair claim is carried by T--N rather than a neutrality assumption.

\paragraph{R$_{\mathrm{surf}}$: exact-output context-materialization control.}
R$_{\mathrm{surf}}$ runs the exact same targeted operation source as T on the same candidate evidence. Rather than exposing the output through \texttt{reusable\_helper\_output}, the harness adds the identical object as \texttt{retrieval\_operation\_output} inside the evidence package and exposes no callable helper. A zero-call parity audit verifies for every tested endpoint that (i) removing this one field recovers the original evidence package exactly, (ii) operation-output content hashes are identical between T and R$_{\mathrm{surf}}$, and (iii) no candidate evidence is removed or reordered. On all 18 pre-frozen DeepSeek load-bearing endpoints, two fresh repeats per arm give T--R$_{\mathrm{surf}}=0.0$ points with bootstrap 90\% CI $[-5.6,5.6]$ and post-hoc 95\% sensitivity CI $[-8.3,8.3]$. A paired-endpoint TOST at the prospectively frozen $\pm10$-point margin gives $p=0.0121$ on each one-sided test; paired-t 95\% CI is $[-8.5,8.5]$. A deterministic parametric sensitivity calculation using the observed endpoint SD (17.15 points) gives about 53\% probability of declaring equivalence when the true mean is exactly zero, so we do not characterize the design as high-powered. C3 grounding is T=R$_{\mathrm{surf}}$=60\% with five endpoint ties; C4 grounding is 70\%/70\% with one win, three ties, and one loss; EIA is ceiling at 100\%/100\%. The four strictly non-ceiling fresh endpoints have mean residual 0 but a wide 90\% interval $[-25,25]$. A secondary Kimi alignment check also ties on all runs. We therefore describe the observed DeepSeek portfolio as consistent with equivalence at the $\pm10$-point \emph{portfolio-average} margin under this forced one-answer harness, with low sensitivity and a ceiling-heavy composition. Fine-grained non-ceiling effects, adaptive callable use, and independently designed temporal retrievers remain outside the claim.

\subsection{R$_{\mathrm{surf}}$ endpoint residuals and future resolution planning}\label{sec:future-power}
The $\pm10$-point threshold was frozen before R$_{\mathrm{surf}}$ outcomes as a materiality threshold and was not derived from the observed T--R$_{\mathrm{surf}}$ variance. The pre-outcome contract does not record a variance-based power rationale. Post review, we note only as contextual scale information that targeted contrasts already visible before the R$_{\mathrm{surf}}$ run were on a 20--60 point scale; this does not retroactively turn that scale argument into a preregistered justification or make the completed $n=18$ design high-powered. Using the observed pooled endpoint SD of 17.15 points, a deterministic post-review planning simulation gives about 53\% probability of passing the same TOST when the true mean is zero at $n=18$, and about 81\% at $n=27$. The latter is a future-design estimate under the same variance, not retroactive evidence. Family-specific variance estimates are too unstable to justify smaller targets: C3 grounding and EIA have zero observed residual variance because all endpoint means tie (EIA also ceilings), whereas the five C4 grounding endpoints have SD 35.36 points; under that observed variance, even $n=100$ gives only about 75\% simulated sensitivity and roughly $n=110$ is needed to reach 81\%. Future non-ceiling work should therefore pre-register an independent endpoint count and stopping rule rather than expand until an equivalence test passes.

\begin{table}[ht]
\centering
\scriptsize
\setlength{\tabcolsep}{3pt}
\begin{tabular}{llllr}
\toprule
Endpoint suffix & Stratum & T (\%) & R$_{\mathrm{surf}}$ (\%) & $\Delta$ (pp) \\
\midrule
128dc1bd & C3 Grd & 100 & 100 & 0 \\
15ffceb0 & C4 Grd & 50 & 0 & +50 \\
187e543f & C4 Grd & 0 & 50 & $-50$ \\
1dfccd3a & C3 Grd & 0 & 0 & 0 \\
5c43ed3f & C4 Grd & 100 & 100 & 0 \\
65c03428 & C3 Grd & 100 & 100 & 0 \\
70e1e3d1 & C3 Grd & 0 & 0 & 0 \\
71b83aed & C4 Grd & 100 & 100 & 0 \\
95037898 & C4 Grd & 100 & 100 & 0 \\
cc293f19 & C3 Grd & 100 & 100 & 0 \\
229defbd & EIA Cut & 100 & 100 & 0 \\
4054808c & EIA Cut & 100 & 100 & 0 \\
62014739 & EIA Cut & 100 & 100 & 0 \\
9bf3bfa4 & EIA Cut & 100 & 100 & 0 \\
a4f03bdc & EIA Cut & 100 & 100 & 0 \\
b246cd01 & EIA Cut & 100 & 100 & 0 \\
c0799777 & EIA Cut & 100 & 100 & 0 \\
d6a5cf91 & EIA Cut & 100 & 100 & 0 \\
\bottomrule
\end{tabular}
\caption{All 18 pre-frozen DeepSeek endpoint-mean R$_{\mathrm{surf}}$ residuals. Endpoint suffixes are shortened display identifiers; full IDs and repeat-level rows are in the supplement.}
\label{tab:rsurf-residuals}
\end{table}

The two non-ties do not point in the same direction. On endpoint \texttt{15ffceb0}, only T repeat 0 selects the registered successful span pair; T repeat 1 and both R$_{\mathrm{surf}}$ repeats select the same failing pair. On \texttt{187e543f}, only R$_{\mathrm{surf}}$ repeat 0 succeeds; both T repeats and R$_{\mathrm{surf}}$ repeat 1 select the same failing pair. The targeted operation-output object is content-identical across T and R$_{\mathrm{surf}}$ for both endpoints. The opposite-signed, single-repeat deviations therefore do not isolate a persistent one-direction surface mechanism.

\subsection{Prospective extension to adaptive multi-turn surfaces}
The current R$_{\mathrm{surf}}$ result is intentionally limited to a forced one-answer harness. A multi-turn extension would be a new experiment, not an inference from the present rows. A valid design should (i) use independent episodes, not turns, as scientific units; (ii) freeze the candidate-evidence stream and operation implementation before outcomes; (iii) separate an \emph{availability} estimand (callable interface present versus absent) from a \emph{placement} estimand in which the exact same operation output is synchronized across callable and context-materialized arms; and (iv) pre-register how endogenous invocation opportunities are aligned, rather than selecting turns after observing successful calls. A same-surface no-op arm can diagnose interface perturbation at each registered opportunity. This would test whether adaptive callable state matters when invocation and context evolve across turns; the present study makes no such claim.

A benign mechanism-agnostic generic helper would likewise be a new treatment. It could be useful for triangulating typical T-versus-generic practice, but it is not retroactively inserted into the current identification ladder: G remains a stress comparator and G$_0$ a same-surface no-operation placebo. Any future benign-generic study must freeze its behavior, endpoint units, and decision rule before new outcomes.

\subsection{Follow-up extension: benign organization, cross-domain replay, and actual tool use}\label{sec:followup-extension}
After the primary audit and its claim boundary were frozen, we ran a separate follow-up extension. These outcomes were \emph{not} used to choose the primary claim and are not retroactively described as preregistered evidence for it. Their role is narrower: test whether a more realistic mechanism-agnostic helper absorbs targeted credit, obtain fresh non-ceiling and cross-domain endpoints, and check whether real multi-turn function availability creates a residual that the forced one-answer R$_{\mathrm{surf}}$ experiment could not measure. All extension calls use the same provider route and exact DeepSeek model identity as the primary fresh controls; each completed call or turn is checkpointed before the next unit. Repeats are averaged within endpoint and never increase the scientific sample size.

\paragraph{A benign generic evidence organizer.}
We introduce a new arm B that is deliberately useful but mechanism-agnostic. B copies and organizes the complete supplied evidence in its original order and adds a generic verification reminder. It is forbidden from filtering by release date, inferring reporting periods, selecting mechanism-specific grounding spans, or reordering candidates. Thus B differs from the historical mechanism-misaligned stress helper G and from the information-empty placebo G$_0$. An initial four-endpoint historical pilot was all ties and therefore not scaled. We then froze genuinely new EIA releases and a new BLS CPI substrate before their respective model outcomes. A first-party EIA archive page dated 2026-06-24 contains an internally inconsistent release-date header; that source anomaly was excluded before model outcomes rather than silently corrected. The clean EIA targets are the July 22, July 29, August 5, and August 12 releases, with the next release used only as a post-cutoff distractor. The BLS set uses four CPI release targets from April through July 2026 and first-party archive-header fields matching the information shown to the model.

\begin{table}[ht]
\centering
\scriptsize
\setlength{\tabcolsep}{4pt}
\begin{tabular}{lrrrrrr}
\toprule
Fresh substrate & N & B & T & T--N & B--N & T--B \\
\midrule
EIA WPSR ($n=4$) & .375 & 1.000 & 1.000 & +.625 & +.625 & .000 \\
BLS CPI ($n=4$) & .375 & .750 & 1.000 & +.625 & +.375 & +.250 \\
\bottomrule
\end{tabular}
\caption{Follow-up extension using endpoint-mean family success after two execution repeats per endpoint. Repeats measure robustness and are averaged within the four independent endpoints. B is the mechanism-agnostic evidence organizer.}
\label{tab:benign-extension}
\end{table}

The EIA result is especially diagnostic: T still repairs the original agent by 62.5 points on average, but B produces the same endpoint means as T on all four endpoints. On this substrate the targeted operation therefore receives no residual credit over generic evidence organization. BLS is mixed rather than uniformly absorbed: generic organization explains part of the net repair, while one target has stable T$>$B=N across both repeats. The two fresh sets therefore sharpen, rather than overturn, the paper's attribution message: T--N can establish net repair without establishing that the registered temporal operation uniquely deserves that credit. Because each set has only four endpoints, these follow-ups are robustness and boundary evidence rather than high-resolution effect estimates.

\paragraph{Actual multi-turn function availability.}
We separately verified that the provider supports a genuine function-call round trip: the model emits a function call, receives a \texttt{function\_call\_output}, and then produces a final answer. On the four clean EIA endpoints, each experimental unit begins from one shared first-turn planning response and then forks into BASE (no helper), CONTEXT (the exact frozen T1 output materialized as ordinary context), or TOOL (a real \texttt{temporal\_filter} function is available and invocation is endogenous). All three arms succeed on 4/4 endpoints in this exploratory fork. TOOL is actually invoked on 2/4 endpoints, yet TOOL--CONTEXT is zero on every endpoint. This is a null boundary rather than evidence for multi-turn equivalence: the planning turn itself changes the baseline, the design is small, and endogenous invocation is part of the TOOL treatment, so TOOL--CONTEXT is not a pure placement estimand.

\paragraph{A planning-baseline hypothesis fails prospectively.}
The unexpectedly strong helper-free BASE arm suggested a possible spin-off hypothesis: perhaps an extra structured planning/decomposition turn absorbs apparent temporal-skill gains. On the EIA endpoints, two planning-only executions give endpoint-mean success .875 versus .375 for the matched single-turn N runs; on the BLS endpoints the corresponding exploratory values are 1.000 versus .375. These observations motivated, but do not confirm, the hypothesis: the EIA result generated the idea, and the BLS planning arm was added after the original BLS outcomes were visible.

We therefore froze a third-substrate falsifier \emph{before any Federal Reserve model outcome}. Four FOMC statement targets (December 10, 2025; January 28, March 18, and April 29, 2026) were selected from one consecutive first-party statement chain, each with three prior statements and one next-statement distractor. The comparison was fixed as single-turn N versus helper-free two-turn planning before either arm ran. Single-turn N succeeds on 1/4 endpoints and planning-only also succeeds on 1/4. Pairwise, planning has one win, N has one win, and two endpoints tie, for mean difference zero. Inspection shows no stable cutoff mechanism: on one endpoint planning removes the future distractor, while on another it introduces one that the single-turn answer had excluded. We therefore \emph{stop} the proposed planning/deliberation paper direction at this point. The positive EIA/BLS observations remain visible exploratory evidence, but they do not survive the prospectively frozen independent-substrate falsifier. Reopening that scientific object would require a new prospectively frozen substrate showing a stable planning-over-single-turn residual under the same-information contract.

\paragraph{Checkpointed execution and recovery.}
Every fresh G$_0$ and R model call is persisted before the next unit begins: an atomic raw JSON receipt, one fsynced CSV row, one fsynced JSONL row, and an atomic checkpoint summary keyed by endpoint$\times$repeat$\times$arm. Resume logic executes missing units only and never reruns a successfully checkpointed unit. Each experiment first runs a small runtime/protocol pilot whose promotion rule ignores scientific outcomes. This design was exercised in practice: transient control-channel failures occurred during the Kimi run while the remote process continued, and the already completed rows remained recoverable from CSV/raw checkpoints without replaying prior calls.

Every retained row binds endpoint, condition, repeat, skill hash, prompt hash, raw response hash, and provider response-ID hash. Kimi repeat 4 is excluded by hash provenance with no replacement. The exact TimeSage-EV replication remains separate from all source-native evidence.
